\section{Realization by an ordinary metric}
\label{sec:correction}

We now turn a sufficiently accurate real quasimode into an exact
eigenfunction of an ordinary four-dimensional metric.  The construction
preserves signs and retains a region on which both the metric and the
eigenfunction are the explicit background ones.  In this section all
\(C^r\) norms are taken in fixed finite coordinate charts.  Constants
may depend on the fixed background and patches, but not on \(k\).

\begin{lemma}[Ordinary-metric realization]\label{corr:realization}
Fix the background \(M,g_0,P,u_k^0,\Lambda_k\) of
Lemma~\ref{env:background-mode}, patches
\(U_2\Subset U_3\Subset U_4\Subset P\), a smooth metric \(h\) such that
\(h-g_0\) has compact support in \(U_2\), and a fixed smooth function
\(f\) on a neighborhood of \(\overline P\).  Put
\[
 W=e^{kf}+e^{kf_0}.
\]
For every integer \(p\ge0\) and every \(L>0\), there are finite
integers \(\ell,J\) with the following property.  Suppose smooth real
functions \(u=u_k\), for arbitrarily large integers \(k\), satisfy
\begin{align}
 &\supp(u-u_k^0)\Subset U_2,
 \qquad \supp R\Subset U_2,
 \qquad R=(\Delta_h+\Lambda_k)u, \label{corr:support}\tag{C1}\\
 &W^{-1}|\partial^\alpha u|
       \le C_\alpha k^{|\alpha|+5}
       \quad\hbox{on }P,\quad |\alpha|\le\ell+1,
       \label{corr:derivatives}\tag{C2}\\
 &\bigl(u^2+k^{-2}|du|_{\mathrm{Euc}}^2\bigr)^{1/2}
       \ge k^{-200}W
       \quad\hbox{on }\overline U_2,\label{corr:noncancellation}\tag{C3}\\
 &W^{-1}|\partial^\alpha R|
       \le C_R k^{-J}
       \quad\hbox{on }P,\quad |\alpha|\le\ell.
       \label{corr:residual}\tag{C4}
\end{align}
The constants in these estimates are independent of \(k\); the supports
in \eqref{corr:support} need not be independent of \(k\).
For all sufficiently large such \(k\), there are a smooth ordinary
metric \(g_*\) and a smooth positive function \(b_*\) on \(M\) such that
\[
 \|g_*-h\|_{C^p}=O(k^{-L}),\qquad
 t=b_*u,\qquad (\Delta_{g_*}+\Lambda_k)t=0.
\]
Moreover, \(g_*-g_0\) and \(t-u_k^0\) have compact support in \(U_4\).
In particular, the zero set and every strict sign comparison of \(u\)
are preserved.
\end{lemma}

The preceding sections supply all the hypotheses of
Lemma~\ref{corr:realization}.  Its proof below does not use the strict
envelope inequality; that inequality was needed to construct the
quasimode.  The essential geometric input here is the clean collar
\(U_4\setminus\overline U_2\), where \(h=g_0\), \(u=u_k^0\), and the
background gradient has a quantitative lower bound.

\subsection{Exact correction with an independent density}

We first allow the principal tensor and divergence density to vary
independently.  The scalar \(\vartheta\) below is a density increment;
the symbol \(d\) continues to denote differentiation.

\begin{lemma}[Weighted correction]\label{corr:weighted}
Under \eqref{corr:support}--\eqref{corr:residual}, there are a smooth
symmetric contravariant tensor \(F\) and a smooth scalar
\(\vartheta\), both compactly supported in \(U_2\), such that
\begin{equation}\label{corr:weighted-equation}
 \divg_h\bigl((h^{-1}+F)\,du\bigr)
       +\Lambda_k(1+\vartheta)u=0.
\end{equation}
For each integer \(0\le r\le\ell-1\), their estimates through order \(r\)
use only derivatives of \(u\) through order \(r+2\) and of \(R\)
through order \(r+1\).  A fixed finite exponent \(E_r\), independent
of \(J\), satisfies
\begin{equation}\label{corr:weighted-estimate}
 \|F\|_{C^r}+\|\vartheta\|_{C^r}
       \le C_r k^{E_r-J};
 \qquad E_r=1000(r+2)\quad\hbox{is sufficient}.
\end{equation}
\end{lemma}

\begin{proof}
Choose a fixed smooth function \(\gamma_0:\mathbb R\to[0,1]\), equal
to one on \([-1/4,1/4]\) and supported in \((-1/2,1/2)\), and set on
\(U_2\)
\[
 \gamma=\gamma_0(k^{200}u/W).
\]
On the support of \(\gamma\),
\eqref{corr:noncancellation} gives
\(|du|_h\ge c k^{-199}W\).  Wherever \(1-\gamma\ne0\),
\(|u|\ge\tfrac14 k^{-200}W\).  Define on \(U_2\)
\begin{align}
 F&=-u\gamma R|du|_h^{-2}h^{-1},\label{corr:F}\\
 \vartheta&=\Lambda_k^{-1}\left(
       -(1-\gamma)\frac R u
       +\divg_h\bigl(\gamma R|du|_h^{-2}\grad_hu\bigr)
                       \right).\label{corr:theta}
\end{align}
Each quotient is used only where its denominator is nonzero.  Extend
the cutoff quotient by zero on the unused part of its domain.  This is
smooth: near a zero of \(du\), the first cutoff is identically zero,
and near a zero of \(u\), the second cutoff is identically zero.
The support condition on \(R\) then permits smooth extension of both
corrections by zero to \(M\), including the divergence term.

Put \(q=\gamma R|du|_h^{-2}\).  The product rule gives the exact
identities
\begin{align*}
 \divg_h(Fdu)&=-\gamma R-u\divg_h(q\grad_hu),\\
 \Lambda_k\vartheta u&=-(1-\gamma)R
                         +u\divg_h(q\grad_hu).
\end{align*}
Adding these identities to the definition of \(R\) proves
\eqref{corr:weighted-equation}, also at the zeros of \(u\), by
smoothness.

For the estimates, use normalized quantities
\[
 a_0=u/W,\qquad \beta=du/W,\qquad r_0=R/W.
\]
All exponentially varying factors cancel from the formulas:
\begin{align}
 F&=-a_0\gamma r_0|\beta|_h^{-2}h^{-1},\label{corr:normalized-F}\\
 q\grad_hu&=\gamma r_0|\beta|_h^{-2}\beta^{\sharp_h},
 \qquad (1-\gamma)R/u=(1-\gamma)r_0/a_0.
 \label{corr:normalized-quotients}
\end{align}
Since \(f,f_0\) are fixed, differentiation of their exponentials and
the reciprocal of their positive sum gives
\[
 |\partial^jW|\le C_j k^jW,
 \qquad |\partial^j(W^{-1})|\le C_j k^jW^{-1}.
\]
Consequently, for every required order \(j\),
\begin{equation}\label{corr:normalized-bounds}
 |\partial^j a_0|\le C_jk^{j+5},\qquad
 |\partial^j\beta|\le C_jk^{j+6},\qquad
 |\partial^j r_0|\le C_jk^{j-J}.
\end{equation}
Here and below a bound on \(\partial^j\) means the bound for every
multi-index of that length.  The chain and quotient rules give, on the
respective cutoff domains,
\begin{align*}
 |\partial^j\gamma|&\le C_jk^{206j},\quad j\ge1,\\
 |\partial^j(a_0^{-1})|&\le C_jk^{200+206j},\\
 |\partial^j(|\beta|_h^{-2})|&\le C_jk^{398+411j}.
\end{align*}
For example, \(|\beta|_h^2\ge c k^{-398}\) on the first domain,
and \(|\partial^j(|\beta|_h^2)|\le C_jk^{j+12}\).
Each term in its differentiated reciprocal has at most \(j\)
differentiated factors, which proves the last displayed bound.
These deliberately generous powers also bound derivatives at cutoff
interfaces.  Substitution into
\eqref{corr:normalized-F}--\eqref{corr:normalized-quotients}, with one
extra derivative for the divergence in \eqref{corr:theta}, proves
\eqref{corr:weighted-estimate}.  In particular, the loss depends only
on the requested derivative order and the fixed exponent in
\eqref{corr:noncancellation}, and does not grow when \(J\) is increased.
\end{proof}

Take the residual accuracy high enough that \(F\) and \(\vartheta\)
are small.  Then \(h^{-1}+F\) is positive definite and
\(1+\vartheta>0\).  Define an ordinary metric \(h_1\), a positive
weight \(w\), a real function \(v\), and a potential \(V_0\) by
\begin{equation}\label{corr:conjugation-data}
 \begin{aligned}
 h_1^{-1}&=\frac{h^{-1}+F}{1+\vartheta},&
 w&=(1+\vartheta)\frac{d\mathrm{vol}_h}{d\mathrm{vol}_{h_1}},\\
 v&=w^{1/2}u,&
 V_0&=\frac{\Delta_{h_1}(w^{1/2})}{w^{1/2}}.
 \end{aligned}
\end{equation}
The quotient of volume elements means the quotient of their positive
densities.  Dividing \eqref{corr:weighted-equation} by its density
yields
\[
 \Delta_{h_1}u+\langle d\log w,du\rangle_{h_1}
                    +\Lambda_ku=0.
\]
If \(s_0=w^{1/2}\), the product rule cancels the first-order term:
\[
 \Delta_{h_1}(s_0u)
   =s_0\Delta_{h_1}u+2\langle ds_0,du\rangle_{h_1}
                         +u\Delta_{h_1}s_0
   =(V_0-\Lambda_k)s_0u.
\]
Thus
\begin{equation}\label{corr:potential-equation}
 \Delta_{h_1}v=(V_0-\Lambda_k)v.
\end{equation}
The quantities \(h_1-h,w-1,V_0\) have compact support in \(U_2\),
and \(v=u_k^0\) outside a compact subset of \(U_2\).
For any fixed integer \(m\), smooth matrix inversion, the determinant
formula for volume, and two differentiations in \(V_0\) give
\begin{equation}\label{corr:potential-estimate}
 \|h_1-h\|_{C^{m+2}}+\|w-1\|_{C^{m+2}}+\|V_0\|_{C^m}
 \le C_m\bigl(\|F\|_{C^{m+2}}+
                         \|\vartheta\|_{C^{m+2}}\bigr),
\end{equation}
whenever the right-hand side is sufficiently small.

\subsection{A conformal equation with a coercive linearization}

For a positive function \(Z\), the four-dimensional conformal
transformation \(g=Z^2h_1\), \(t=v/Z\), gives
\begin{equation}\label{corr:conformal-identity}
 \Delta_gt
   =Z^{-4}\divg_{h_1}\bigl(Z^2\grad_{h_1}(v/Z)\bigr)
   =Z^{-4}\bigl(Z\Delta_{h_1}v-v\Delta_{h_1}Z\bigr).
\end{equation}
Together with \eqref{corr:potential-equation}, this shows that the
equation
\begin{equation}\label{corr:Z-equation}
 \Delta_{h_1}Z=V_0Z+\Lambda_k(Z^3-Z)
\end{equation}
is sufficient for \((\Delta_g+\Lambda_k)t=0\).

\begin{lemma}[The conformal factor]\label{corr:conformal}
Fix the smooth metric \(h\) on the closed four-manifold \(M\), and
an even integer \(s>2\).  There are \(\delta>0\) and \(C_s\),
depending only on \(h,s\) and the fixed Sobolev norms, with the
following property.  If \(h_1,V_0\) are smooth,
\[
 \|h_1-h\|_{C^s}+\|V_0\|_{H^s}\le\delta,
 \qquad \Lambda\ge1,
\]
then
\[
 \Delta_{h_1}Z=V_0Z+\Lambda(Z^3-Z)
\]
has a smooth positive solution satisfying
\begin{equation}\label{corr:Z-smallness}
 \|Z-1\|_{H^s}\le C_s\Lambda^{-1}\|V_0\|_{H^s}.
\end{equation}
\end{lemma}

\begin{proof}
Use the real Sobolev spaces of the fixed metric \(h\).  The standard
elliptic spectral and Sobolev theorems give the eigenbasis for
\(-\Delta_h\), the equivalence of the usual even-order Sobolev norms
with the norms defined by powers of \(1-\Delta_h\), and the algebra
property of \(H^s\) for \(s>2\)
\cite{Taylor2011,TaylorPsido}.  If \(\mu\ge0\) is a spectral value of
\(-\Delta_h\), the inverse of
\(L_\Lambda=\Delta_h-2\Lambda\) has multiplier
\(-1/(\mu+2\Lambda)\).  In the spectral Sobolev norms, therefore,
\begin{align}
 \|L_\Lambda^{-1}\|_{H^s\to H^s}
       &\le (2\Lambda)^{-1},\label{corr:mass-inverse-one}\\
 \|L_\Lambda^{-1}\|_{H^{s-2}\to H^s}
       &\le\sup_{\mu\ge0}\frac{1+\mu}{\mu+2\Lambda}
        \le1.\label{corr:mass-inverse-two}
\end{align}
Changing to the fixed coordinate Sobolev norms costs constants
independent of \(\Lambda\).

Writing \(Z=1+Y\), the required fixed-point equation is
\begin{equation}\label{corr:fixed-point}
 Y=L_\Lambda^{-1}\left[
    (\Delta_h-\Delta_{h_1})Y+V_0(1+Y)
                            +\Lambda(3Y^2+Y^3)\right].
\end{equation}
The coefficient formula for a Laplacian and the usual integer-order
product rule show
\[
 \|\Delta_h-\Delta_{h_1}\|_{H^s\to H^{s-2}}
                         \le C_s\|h_1-h\|_{C^s}.
\]
Estimate this term using \eqref{corr:mass-inverse-two}.  Estimate the
remaining terms using the algebra property and
\eqref{corr:mass-inverse-one}.  On the closed \(H^s\)-ball of radius
\(r\), the Lipschitz constant of the right side of
\eqref{corr:fixed-point} is at most
\[
 C_s\|h_1-h\|_{C^s}
   +C_s\Lambda^{-1}\|V_0\|_{H^s}+C_s(r+r^2).
\]
Choose a fixed small \(r\) so that the last term is at most \(1/4\)
and Sobolev embedding makes \(1+Y\ge1/2\) on this ball.  Next choose
\(\delta\) so that the displayed Lipschitz constant is at most
\(1/2\) and the image of zero has norm at most \(r/2\).  The map is
then a contraction of the ball into itself.  Its fixed point obeys
\[
 \|Y\|_{H^s}
 \le\tfrac12\|Y\|_{H^s}
                 +C_s\Lambda^{-1}\|V_0\|_{H^s},
\]
which proves \eqref{corr:Z-smallness}.  Finally the equation for
\(Z\), the smoothness of its coefficients, and elliptic regularity
bootstrap \(Z\) from \(H^s\) to \(H^{s+2}\), and then to all higher
orders.  The resulting real smooth solution remains positive.
\end{proof}

The factor \(\Lambda\) in the nonlinear term is canceled by the first
inverse bound, while the second inverse bound handles the change of
the principal coefficients.  These are estimates for the negative
operator \(\Delta_h-2\Lambda\); no high-energy nonresonance assumption
is involved.

\subsection{Localization and exact determinant compatibility}

The global factor \(Z\) need not equal one outside \(P\).  We localize
it, keeping the resulting equation exact by recording its density.
Choose a fixed \(\chi\in C_c^\infty(U_4;[0,1])\) equal to one on a
neighborhood of \(\overline U_3\), and set
\begin{equation}\label{corr:localized-data}
 \widehat Z=1+\chi(Z-1),\qquad
 t=\frac v{\widehat Z},\qquad
 \rho=\widehat Z^2\left(1-\frac{V_0}{\Lambda_k}\right)
                +\frac{\widehat Z\Delta_{h_1}\widehat Z}{\Lambda_k}.
\end{equation}
For sufficiently small \(Z-1\) and \(V_0\), both \(\widehat Z\)
and \(\rho\) are positive.  Direct differentiation gives
\begin{equation}\label{corr:localized-flux}
 \divg_{h_1}(\widehat Z^2h_1^{-1}dt)
       =\widehat Z\Delta_{h_1}v-v\Delta_{h_1}\widehat Z
       =-\Lambda_k\rho t.
\end{equation}
Where \(\chi\) is identically one, \eqref{corr:Z-equation} implies
\(\rho=\widehat Z^4\).  Outside \(\supp\chi\), the same equality
holds because \(V_0=0\) and \(\widehat Z=1\).  Thus
\begin{equation}\label{corr:defect-support}
 \supp(\rho-\widehat Z^4)
        \Subset \mathcal C:=U_4\setminus\overline U_2.
\end{equation}
On this entire collar \(h_1=g_0\) and \(v=u_k^0\).

\begin{lemma}[Determinant repair on the collar]\label{corr:localization}
For the data \eqref{corr:localized-data}, suppose that
\(\widehat Z-1\) is sufficiently small in \(C^1\) and
\(\rho>0\).  Then an ordinary smooth metric \(g_*\), equal to
\(g_0\) outside a compact subset of \(U_4\), satisfies
\((\Delta_{g_*}+\Lambda_k)t=0\).
More quantitatively, if for some \(N>0\),
\begin{equation}\label{corr:collar-input-estimate}
 \|h_1-h\|_{C^p}+\|V_0\|_{C^p}=O(k^{-N}),
 \qquad \|Z-1\|_{C^{p+2}}=O(k^{-N-2}),
\end{equation}
then
\begin{equation}\label{corr:final-metric-estimate}
 \|g_*-h\|_{C^p}
             =O(k^{-N}+k^{p-N-2}).
\end{equation}
\end{lemma}

\begin{proof}
Put \(\phi=\theta+z\).  Since
\(u_k^0=e^{kf_0}\cos(k\phi)\), the normalized covector on
\(\mathcal C\) is exactly
\begin{equation}\label{corr:normalized-collar-covector}
 \alpha:=\frac{dt}{ke^{kf_0}}
  =\widehat Z^{-1}
       \bigl(\cos(k\phi)\,df_0-\sin(k\phi)\,d\phi\bigr)
   -k^{-1}\cos(k\phi)\widehat Z^{-2}d\widehat Z.
\end{equation}
On \(\overline U_4\), the background covectors \(df_0\) and
\(d\phi\) are orthogonal in \(g_0\), with lengths bounded below.
It follows that \(|\alpha|_{h_1}\ge c>0\) for all sufficiently large
\(k\).  In particular,
\begin{equation}\label{corr:clean-gradient}
 |dt|_{h_1}\ge ck e^{kf_0}\quad\hbox{on }\mathcal C.
\end{equation}
Let
\[
 \Pi=I-\frac{\alpha^{\sharp_{h_1}}\otimes\alpha}
                       {|\alpha|_{h_1}^2}
\]
be the \(h_1\)-orthogonal projection of the tangent bundle onto
\(\ker dt\) on \(\mathcal C\).  It has rank three and annihilates
\(\grad_{h_1}t\).

Define the positive scalar
\[
 c_* =\left(\frac{\rho^2}{\widehat Z^8}\right)^{1/3}.
\]
By \eqref{corr:defect-support}, \(c_*-1\) has compact support in
\(\mathcal C\).  Consequently \((c_*-1)\Pi\) extends smoothly
by zero to \(M\).  Define a positive symmetric contravariant tensor
and an ordinary metric by
\begin{equation}\label{corr:determinant-tensor}
 A'=\widehat Z^2\bigl(I+(c_*-1)\Pi\bigr)h_1^{-1},
 \qquad g_*^{-1}=\rho^{-1}A'.
\end{equation}
Here the product containing \(\Pi\) always means this smooth zero
extension.  In the normal direction to \(\ker dt\), the endomorphism
\(A'h_1\) has eigenvalue \(\widehat Z^2\); in the three tangent
directions its eigenvalue is \(\widehat Z^2c_*\).  Hence
\begin{equation}\label{corr:flux-and-determinant}
 A'dt=\widehat Z^2h_1^{-1}dt,
 \qquad \det(A'h_1)=\widehat Z^8c_*^3=\rho^2.
\end{equation}
Off the collar the same identities follow from \(c_*=1\) and
\(\rho=\widehat Z^4\).  Dimension four now gives
\[
 \det(g_*^{-1}h_1)=\rho^{-4}\det(A'h_1)=\rho^{-2},
 \qquad d\mathrm{vol}_{g_*}=\rho\,d\mathrm{vol}_{h_1}.
\]
It follows from \eqref{corr:localized-flux} and
\eqref{corr:flux-and-determinant} that
\[
 \Delta_{g_*}t
   =\rho^{-1}\divg_{h_1}(A'dt)=-\Lambda_kt.
\]
All changes are compactly supported in \(U_4\).

It remains to verify the estimates without losing an exponential
factor from \(dt\).  Under \eqref{corr:collar-input-estimate},
\(\widehat Z\ge1/2\) and its derivatives through order \(p+2\)
are bounded independently of \(k\).  Formula
\eqref{corr:normalized-collar-covector} gives
\[
 |\partial^j\alpha|\le C_jk^j,
 \qquad |\partial^j\Pi|\le C_jk^j,
 \qquad 0\le j\le p.
\]
The second estimate follows from the first, the fixed positive lower
bound on \(|\alpha|\), and the quotient rule.  These powers are
independent of the accuracy requested in the residual.
On the region where \(\chi\equiv1\), use \(\rho=Z^4\).  In the
transition region use \eqref{corr:localized-data} with the fixed
coefficients \(h_1=g_0\), \(V_0=0\), and outside \(\supp\chi\)
use \(\rho=1\).  Since \(\Lambda_k\) is comparable to \(k^2\),
\eqref{corr:collar-input-estimate} therefore gives
\[
 \|\rho-1\|_{C^p}+\|c_*-1\|_{C^p}=O(k^{-N-2}).
\]
The product rule therefore yields
\[
 \|(c_*-1)\Pi\|_{C^p}=O(k^{p-N-2}).
\]
For the covariant metric, use the exact inverse supplied by
\(\Pi^2=\Pi\):
\[
 g_*=\rho\widehat Z^{-2}h_1
                  \bigl(I+(c_*^{-1}-1)\Pi\bigr).
\]
The scalar \(c_*^{-1}-1\), like \(c_*-1\), has \(C^p\) norm
\(O(k^{-N-2})\).  The product rule in this formula, together with
\(\|h_1-h\|_{C^p}=O(k^{-N})\), proves
\eqref{corr:final-metric-estimate} for every \(N>0\), including when
its second term grows with \(k\).
\end{proof}

\subsection{The finite order of choices}

\begin{proof}[Proof of Lemma~\ref{corr:realization}]
Here is one explicit order of choices.  Given \(p,L\), choose an even
integer \(s>p+4\), and set
\[
 \begin{gathered}
 N=\lceil L\rceil+p+4,\qquad r=s+2,\qquad \ell=s+3,\\
 J=E_{s+2}+N,\qquad E_{s+2}=1000(s+4).
 \end{gathered}
\]
Lemma~\ref{corr:weighted} uses \(u\) through order
\(s+4=\ell+1\) and \(R\) through order \(\ell\), and gives
\[
 \|F\|_{C^{s+2}}+\|\vartheta\|_{C^{s+2}}=O(k^{-N}).
\]
Equation~\eqref{corr:potential-estimate} then gives
\[
 \|h_1-h\|_{C^{s+2}}+\|w-1\|_{C^{s+2}}
                           +\|V_0\|_{H^s}=O(k^{-N}).
\]
For large \(k\), Lemma~\ref{corr:conformal} applies with
\(\Lambda=\Lambda_k\).  Its estimate and the embedding
\(H^s\hookrightarrow C^{p+2}\), valid because \(s>p+4\) in
dimension four, give
\[
 \|Z-1\|_{C^{p+2}}
       \le C\|Z-1\|_{H^s}=O(k^{-N-2}).
\]
The hypotheses of Lemma~\ref{corr:localization} now hold.  Since
\(N\ge L+p+4\), its conclusion implies
\(\|g_*-h\|_{C^p}=O(k^{-L})\).  The eigenfunction is exactly
\[
 t=\frac{w^{1/2}}{\widehat Z}\,u,
\]
with a smooth positive multiplier equal to one off a compact subset
of \(U_4\).  This proves every assertion.

All derivative orders and powers have been fixed before the wave
construction is asked to provide
\eqref{corr:derivatives}--\eqref{corr:residual}.  The constants and
the eventual lower threshold for \(k\) may depend on that finite
construction.  Increasing its Taylor orders changes those constants,
but neither the fixed noncancellation exponent nor the polynomial
losses above.  Thus only finitely many accuracy requirements are used
at each realization step.
\end{proof}
